\newpage
\unnumberedtitle{RESUMO}

\begingroup
\setstretch{1}
\setlength{\parindent}{0pt}

\authorsCitation{} \textbf{\workTitle}\ifdefempty{\workSubtitle}{}{: \workSubtitle}. \workYear. \pageref{LastPage} f. Trabalho de Conclusão de Curso (Bacharelado em \courseName) -- \centerName, \universityName, \workCity, \workYear.\par
\vspace{\baselineskip}

A gestão financeira pessoal no Brasil ainda depende da consulta a diferentes aplicativos e da interpretação de dados dispersos entre várias instituições. Embora o Open Finance permita o compartilhamento padronizado de dados financeiros mediante consentimento do titular, a maior parte das soluções disponíveis concentra-se na visualização estática de informações. O objetivo foi desenvolver o EIGAM, um assistente financeiro inteligente, preditivo e conversacional voltado ao público jovem, capaz de operar dentro das ferramentas de conversa que o usuário já utiliza. A arquitetura separa o núcleo da solução do provedor de mensageria: um adaptador incorpora novos aplicativos de \textit{chat} sem alterar o núcleo, e essa independência se confirmou com dois mensageiros de mercado e com uma conversa no painel \textit{web} do titular. Um servidor Model Context Protocol (MCP) permite, ainda, que o assistente de inteligência artificial do próprio titular consulte os dados. A solução adota uma arquitetura de microsserviços orientada a eventos, com quatorze serviços em TypeScript e Python, e se integra ao Open Finance por meio de um agregador de dados. Um planejador baseado em modelo de linguagem escolhe as consultas financeiras necessárias, consultas estruturadas calculam os valores, regras determinísticas aplicam o mascaramento e o titular exerce por mensagem os direitos previstos na Lei Geral de Proteção de Dados Pessoais (LGPD). Avisos de limite de gasto, de cobranças duplicadas e de transações atípicas, metas com registro reversível, resumos periódicos e projeções de gasto e de saldo livre tornam o acompanhamento proativo, e uma simulação de compra estima, a pedido do titular, o efeito de uma compra no orçamento. A avaliação da qualidade abrangeu testes automatizados com cobertura exigida pela integração contínua, conjuntos de referência rotulados para a camada conversacional, um teste de carga no ambiente de produção e testes retrospectivos das projeções sobre séries sintéticas, dados reais e uma base bancária pública. As regras determinísticas de mascaramento de valores e de reconhecimento dos comandos da LGPD atingiram conformidade total nos conjuntos de referência, e a avaliação das respostas abertas incluiu um conjunto nunca usado para ajustes. O sistema só exibe estimativas aprovadas nos critérios estatísticos definidos; sem aprovação, ele se abstém em vez de apresentar números pouco confiáveis.\par
\vspace{\baselineskip}

\textbf{Palavras-chave:} Open Finance; serviços de mensageria; inteligência artificial; análise preditiva; assistência financeira conversacional.
\endgroup

\newpage
\unnumberedtitle{ABSTRACT}

\begingroup
\setstretch{1}
\setlength{\parindent}{0pt}

\authorsCitation{} \textbf{\workTitleEnglish}\ifdefempty{\workSubtitleEnglish}{}{: \workSubtitleEnglish}. \workYear. \pageref{LastPage} f. Undergraduate Thesis (Bachelor's degree in \courseNameEnglish) -- \centerNameEnglish, \universityNameEnglish, \workCity, \workYear.\par
\vspace{\baselineskip}

Personal financial management in Brazil still depends on checking several applications and interpreting data scattered across different institutions. Although Open Finance enables the standardized sharing of financial data with the account holder's consent, most available solutions focus on the static visualization of information. The objective was to develop EIGAM, an intelligent, predictive and conversational financial assistant aimed at young users, able to operate inside the chat tools users already have. The architecture decouples the solution's core from the messaging provider: an adapter adds new chat applications without changing the core. This independence was confirmed with two mainstream messengers and a conversation in the account holder's web panel. A Model Context Protocol (MCP) server also lets the holder's own AI assistant read their data. The solution adopts an event-driven microservices architecture, with fourteen services written in TypeScript and Python, and integrates with Open Finance through a data aggregator. A language-model-based planner selects the required financial queries, structured queries compute the amounts, deterministic rules apply masking and the account holder exercises the rights granted by the Brazilian General Data Protection Law (LGPD) through messages. Spending limit alerts, duplicate charge and unusual transaction alerts, goals with a reversible ledger, periodic summaries and spending and free-balance forecasts make monitoring proactive, and a purchase simulation estimates, on request, the effect of a purchase on the budget. The quality assessment covered automated tests with coverage enforced by continuous integration, labeled reference sets for the conversational layer, a load test in the production environment and backtests of the forecasts over synthetic series, real data and a public banking dataset. The deterministic rules for amount masking and LGPD command recognition reached full compliance on the reference sets, and the evaluation of open-ended answers included a set never used for tuning. The system only shows estimates that pass the defined statistical criteria; otherwise, it abstains instead of presenting unreliable figures.\par
\vspace{\baselineskip}

\textbf{Keywords:} Open Finance; messaging services; artificial intelligence; predictive analytics; conversational financial assistance.
\endgroup
